\begin{abstract}
Multi-agent pipelines are increasingly deployed to filter the outputs of large
language models (LLMs): a team of LLM agents inspects a candidate response, and a
coordinator combines their verdicts into an allow/block decision. Existing designs,
exemplified by AutoDefense, sharply reduce jailbreak attack success rates, but they
secure the \emph{content} while leaving the \emph{defender itself} unexamined. Such
pipelines assume every internal agent is honest, route all verdicts through a single
trusted coordinator with no quorum or Byzantine tolerance, feed the untrusted
candidate directly to the judges that must evaluate it, are stress-tested only against
static attacks, and target jailbreak (harmful content) rather than prompt injection
(task-integrity violations that can appear benign). We take the position that a defense
pipeline is a multi-agent system whose own security must be analysed, and we study two
threats that have no single-model analogue: compromised or colluding \emph{internal}
defenders, and \emph{second-order} prompt injection carried by the very content the
judges read. We present \Aegis, a defense architecture that (i)~replaces the single
coordinator with Byzantine-robust aggregation of judge verdicts (Krum-style selection,
coordinate-wise median, or geometric median), (ii)~structurally isolates the untrusted
payload from each judge's instruction channel, (iii)~hardens each judge against
injection, and (iv)~generalises the verdict target from harmful-content classification
to instruction-hierarchy / task-integrity. We formalise the semantic verdict space and
prove that, under stated honest-concentration and decision-margin assumptions, robust
aggregation with $2f{+}2<n$ (selection) or $f<n/2$ (median rules) prevents up to $f$
Byzantine judges from flipping the aggregate decision in either direction; we further
characterise how the guarantee degrades under correlated judge failures, making the
required backbone diversity explicit rather than assumed. We give a security-versus-cost
analysis showing token and latency cost linear in the committee size, and an adaptive
evaluation protocol covering compromised judges, colluding minorities, second-order
payload injection, and attacks on the aggregation rule. We are explicit that a reduced
attack success rate under compromise is evidence of graceful degradation, not a proof
of end-to-end security; residual failures are reported. The reported empirical figures
are placeholders that specify the intended protocol and must be replaced with measured
results before submission.
\end{abstract}

\vspace{-0.4em}
\noindent\textbf{Highlights}
\begin{itemize}[leftmargin=1.3em,itemsep=1pt,topsep=2pt]
  \item Frames the multi-agent LLM defense pipeline as an object whose own security,
        under insider compromise and second-order injection, must be analysed.
  \item \Aegis: Byzantine-robust verdict aggregation, payload isolation, and hardened
        judges, generalised from jailbreak to prompt-injection/task-integrity.
  \item Decision-integrity guarantee for robust aggregation in a semantic verdict space
        under explicit honest-concentration and margin assumptions.
  \item Correlated-failure analysis that formalises the backbone-diversity requirement
        and shows added judges do not help once errors are correlated.
  \item Adaptive-attack evaluation protocol and a security-versus-cost frontier over the
        number of hardened defenders.
\end{itemize}
